
\subsubsection{Interpretation of Main Finding}

The H-M2 result (Structured > Scrambled, p < 0.001) provides evidence that organizational structure affects repair success independently of information content. The scrambled condition contains identical diagnostic information; only section order differs. This design controls for the possibility that structured formatting improves performance merely by adding information.

The effect size (Cohen's d = 0.30) is small-to-medium. The practical magnitude---14.4 percentage points---represents a meaningful improvement in repair success rate, though substantial room for improvement remains.

\subsubsection{Scaffolding Pattern}

The H-M3 simulation results are consistent with scaffolding theory from human-computer interaction research. The inverted-U pattern suggests that intermediate hints (Level 2) may activate relevant model knowledge without creating copy-paste dependency. However, because these results are from simulation mode, this finding should be considered preliminary.

\subsubsection{Scale Interaction Non-Finding}

The directional pattern (7B > 34B > GPT-4) is consistent with the hypothesis that smaller models benefit more from structured formatting. However, the interaction does not reach statistical significance. Two explanations are plausible:

\begin{enumerate}
\item \textbf{Ceiling effects}: Larger models may already parse errors effectively, leaving less room for format improvement.
\item \textbf{Power limitation}: GPT-4 had only 81 failure samples due to high base accuracy, reducing statistical power for interaction detection.
\end{enumerate}

The current data cannot distinguish between a true null effect and an effect too small to detect at current sample sizes.

\subsubsection{Limitations}

\textbf{Simulation mode for H-M3}: The fix specificity experiment used synthetic data due to flash\_attn CUDA symbol errors. Real-model validation is required.

\textbf{Mock LLM judge for H-M1}: The reconstruction test used regex-based extraction rather than an actual LLM judge due to missing API credentials. While perfect accuracy in mock mode validates format unambiguity, real LLM validation would strengthen this finding.

\textbf{Single language}: All experiments use Python. Results may not generalize to statically-typed languages with different error message structures.

\textbf{Benchmark-specific}: Evaluation uses HumanEval+/MBPP+ consisting of short, self-contained problems. Production codebases involve longer contexts and multi-file dependencies.

\textbf{Error type coverage}: Eight common Python error types are supported. Rare error types and custom exceptions are not covered.

\textbf{Mode}: The H-M2 experiment ran in mock mode (pipeline validation) rather than with full model inference. While the statistical methodology is validated, results should be confirmed with actual model execution.
